% Source-aligned English for printed p.143.
\SourcePage{143}

\LayerTarget{tohoku:remark:satellites-small-skeleton:p2}
\noindent that are abelian and defined over $k$, and whose connected component
of the identity is an ``abelian variety.'' The characteristic must be zero,
since otherwise a bijective homomorphism need not be an isomorphism. We merely
indicate that the preceding result is proved by constructing $S^1F(A)$ as the
direct limit of the objects $F(M/A)/\mathrm{Im}(F(M))$ over ``all''
monomorphisms $A\longrightarrow M$ in $C$, preordered by declaring
$A\longrightarrow M$ to be ``dominated'' by $A\longrightarrow M'$ when there
is a morphism $M\longrightarrow M'$ inducing the identity on $A$.

\LayerTarget{tohoku:section:2.3}
\noindent\textbf{2.3. Derived functors.} Let $C$ and $C'$ be abelian
categories. The theory of derived functors of an additive functor
$F:C\longrightarrow C'$ develops as in [6, Chap. V], except that, according
to the case, one must suppose that every object $A\in C$ is isomorphic to a
subobject of an injective object, to a quotient object of a projective object,
or both.

\LayerTarget{tohoku:construction:right-derived-functors:p1}
Thus, to define the right derived functors of a covariant functor or the left
derived functors of a contravariant functor, one must suppose that every
$A\in C$ is isomorphic to a subobject of an injective object. It then follows
that every $A\in C$ has an injective resolution\footnote{A \textit{right
resolution} of an object $A\in C$ is, by definition, a complex $C$ in
nonnegative degrees equipped with an ``augmentation homomorphism''
$A\longrightarrow C$, where $A$ is regarded as a complex concentrated in
degree zero, such that
$0\longrightarrow A\longrightarrow C^0\longrightarrow C^1\longrightarrow
\cdots$ is exact. An \textit{injective resolution} of $A$ is a resolution $C$
of $A$ for which all $C^i$ are injective objects. Left resolutions of $A$,
and projective resolutions in particular, are defined dually.}:
$0\longrightarrow A\longrightarrow C^0\longrightarrow C^1\longrightarrow
\cdots$. One defines
$R^iF(A)=H^i(F(C))$, where $C$ denotes the complex of the $C^i$. To define
the left derived functors of a covariant functor or the right derived functors
of a contravariant functor, one must likewise suppose that every $A\in C$ is
isomorphic to a quotient of a projective object. Finally, to define derived
functors of a multivariable functor with mixed variance, the appropriate
hypothesis must be imposed on the category of each variable. Apart from this,
the account in [6] carries over unchanged.

\LayerTarget{tohoku:morphism:satellite-to-derived}
In particular, suppose for example that $F$ is covariant and that its right
derived functors $R^iF$ can be formed. Setting $R^iF=0$ for $i<0$, the system
$RF=(R^iF)$ is a cohomological functor, called the \textit{right derived
cohomological functor of $F$}, and there is a canonical morphism of
$\partial$-functors in nonnegative degrees
$SF\longrightarrow RF$. Here $SF=(S^iF)$ is the universal
$\partial$-functor in nonnegative degrees associated to the satellites of
$F$, whose existence follows from Theorem 2.2.2. This morphism is an
isomorphism if and only if $F$ is left exact. It seems that the $R^iF$ are of
little interest except when $F$ is left exact, that is, when they coincide
with the satellite functors. Nevertheless, their simultaneous definition by
injective resolutions is more manageable than the recursive definition of
the $S^iF$, and is particularly well suited to constructing the most important
spectral sequences; see 2.4.

\LayerTarget{tohoku:construction:derived-bifunctor:p1}
Let $C_1,C_2,C'$ be three abelian categories and let $T(A,B)$ be an additive
bifunctor from $C_1\times C_2$ to $C'$, contravariant in $A$ and covariant in
$B$, to fix ideas. Suppose that every object of $C_2$ is isomorphic to a
subobject of an injective object. One can then construct the

\TohokuContinuation{144}
